\chapter{Introduction}
\label{chap:introduction}

\section{What is ActSim?}

\actsim is a Python-based actuarial simulation package designed for loss
simulation, stochastic risk modelling, synthetic claim generation, and
actuarial validation workflows. The package provides a transparent and
extensible framework for actuaries, researchers, and analysts who need to
generate, analyse, and validate simulated insurance loss data.

The package is intended to modernise and extend traditional actuarial
loss simulation tools by providing a Python-native framework that integrates
with modern actuarial analytics workflows, including \texttt{pandas},
\texttt{NumPy}, \texttt{SciPy}, \texttt{chainladder}, and LLM-assisted
actuarial tools.

\section{Why ActSim Exists}

Actuarial workflows have historically depended on a mix of legacy spreadsheet
models, vendor software, and bespoke scripts. While these approaches remain
valuable, they often suffer from limited reproducibility, opaque assumptions,
and weak integration with the broader scientific Python ecosystem.

\actsim addresses several practical needs:

\begin{itemize}[leftmargin=*]
    \item Provide a Python-native framework for frequency--severity simulation
          that can be inspected, audited, and extended.
    \item Offer a reproducible workflow for synthetic claim-level data
          generation, suitable for testing reserving methods and stress
          scenarios.
    \item Support claim development through user-supplied loss development
          factors, with stochastic variation around expected patterns.
    \item Integrate cleanly with triangle-based reserving tools such as the
          \texttt{chainladder} package, so that simulated data can be
          analysed using established reserving methods.
    \item Provide standard risk metrics such as Value at Risk (VaR),
          Tail Value at Risk (TVaR), Aggregate Exceedance Probability (AEP),
          and Occurrence Exceedance Probability (OEP).
\end{itemize}

\section{Capabilities at a Glance}

\actsim supports the following capabilities:

\begin{itemize}[leftmargin=*]
    \item Frequency--severity simulation using a wide range of standard
          distributions.
    \item Aggregate loss modelling with optional dependency between
          frequency and severity, including linear correlation and copula
          structures.
    \item Multivariate correlated simulation across multiple lines of
          business.
    \item Claim-level synthetic data generation from a policy table,
          including occurrence date assignment via a non-homogeneous
          Poisson process.
    \item Claim development using user-defined loss development factors
          with stochastic volatility.
    \item Risk metric calculation, including VaR, TVaR, AEP, and OEP.
    \item Validation and testing workflows, with reproducibility ensured
          via explicit random seeds.
    \item Integration with triangle-based reserving tools.
\end{itemize}

\section{Positioning Statement}

\actsim is intended to support actuarial simulation workflows, including
frequency--severity modelling, stochastic aggregate loss simulation,
claim-level synthetic data generation, claim development, and integration
with triangle-based reserving workflows. It is designed to be seamlessly 
integrated with existing reserving and pricing tools.

\section{Audience}

This manual is written for:

\begin{itemize}[leftmargin=*]
    \item Practising actuaries building loss models in Python.
    \item Researchers who require reproducible synthetic claim data for
          methodology development.
    \item Students and educators who want to demonstrate stochastic
          actuarial techniques with transparent code.
    \item Open-source contributors and reviewers who wish to inspect or
          extend the package.
\end{itemize}

\section{How to Read This Manual}

The manual is organised so that new users can follow it from the
installation chapter through quick start and core concepts before exploring
the main classes. Experienced users can jump directly to the API reference
or worked examples. The validation chapter is intended to support
peer review and open-source scrutiny.
